% Generated draft: requires mathematical and source review.
\documentclass[11pt]{article}
\usepackage[margin=25mm]{geometry}
\usepackage{fontspec}
\setmainfont{Latin Modern Roman}
\usepackage{amsmath,amssymb,mathtools,unicode-math}
\setmathfont{Latin Modern Math}
\usepackage{enumitem,needspace,xurl,hyperref}
\setlength{\parindent}{0pt}
\setlength{\parskip}{5pt}
\setlength{\emergencystretch}{3em}
\newcommand{\N}{\mathbb{N}}
\newcommand{\Z}{\mathbb{Z}}
\newcommand{\Q}{\mathbb{Q}}
\newcommand{\R}{\mathbb{R}}
\newcommand{\C}{\mathbb{C}}
\newcommand{\F}{\mathbb{F}}
\newcommand{\A}{\mathbb{A}}
\newcommand{\PP}{\mathbb{P}}
\newcommand{\E}{\mathbb{E}}
\newcommand{\eps}{\varepsilon}
\newcommand{\dd}{\,\mathrm{d}}
\begin{document}
\section*{UnsolvedMath ID: 30005445}
\textbf{Source problem code:} catalogue entry
\subsection{Definitions and mathematical statement}
The supplied record asks two related questions about random walks in random environments: whether an invariance principle extends to a slightly perturbed balanced environment, and whether it extends to a balanced environment subjected to a percolation perturbation. The motivating objective is to connect the balanced regime discussed in the source report with other regimes, particularly the perturbative regime. The supplied excerpt does not specify the underlying model or either perturbation, so it does not determine a uniquely quantified mathematical conjecture.

To explain the terminology conditionally, suppose a discrete-time random walk takes values in a countable set \(S\subseteq\R^d\), with environment-dependent transition probabilities \(p_\omega(x,y)\). Whenever the relevant first moment is finite, a local zero-drift balance condition would be
\[
\sum_{y\in S}(y-x)p_\omega(x,y)=0\qquad(x\in S).
\]
If the intended model is nearest-neighbour motion on \(\Z^d\), the coordinatewise balance condition is
\[
p_\omega(x,x+e_i)=p_\omega(x,x-e_i)
\qquad(x\in\Z^d,\ 1\leq i\leq d),
\]
where \(e_i\) denotes the \(i\)-th coordinate vector. These are conditional explanations; the excerpt does not identify the intended state space, transition rules, or definition of balance.

For a specified environment law \(\mathbb P\), let \(P_\omega^x\) denote the walk law in environment \(\omega\), started at \(x\). A typical diffusive invariance-principle target, if this scaling and centering are intended, is
\[
\left(\frac{X_{\lfloor nt\rfloor}-X_0-\lfloor nt\rfloor v}{\sqrt n}\right)_{t\geq0}
\Longrightarrow (B_t)_{t\geq0}
\quad\text{in }D([0,\infty),\R^d),
\]
with the Skorokhod \(J_1\) topology, where \(v\in\R^d\) is a deterministic velocity and \(B\) is a Brownian motion started at zero with covariance \(\mathbb E[B_tB_t^{\mathsf T}]=t\Sigma\). A quenched formulation requires convergence under \(P_\omega^x\) for \(\mathbb P\)-almost every environment; an annealed formulation requires convergence under the averaged law \(\int P_\omega^x\,\mathbb P(\dd\omega)\). The source excerpt does not specify either formulation, the starting-point quantifier, the centering, the topology, or whether \(\Sigma\) must be deterministic and positive definite. Thus this display is an explanatory template, not an assertion of the exact intended theorem.

\textbf{Statement retained from the source.} Determine whether an invariance principle can be established for each of the following extensions of the balanced regime in the report: (i) a slightly perturbed balanced environment; (ii) a balanced environment subjected to a percolation perturbation. A precise theorem requires the report's underlying hypotheses and definitions of both perturbations. In particular, the supplied record gives no perturbation metric, smallness threshold, percolation construction, parameter range, conditioning, or rule for the walk on the perturbed environment. These missing specifications remain unresolved.
\subsection{Short English statement}
The problem asks whether invariance principles extend from the report's balanced regime to small perturbations and to percolation perturbations. The supplied excerpt leaves the model, perturbations, assumptions, and exact convergence statement unspecified, so the displayed mathematical framework is conditional and requires human review.
\subsection{Sources}
\begin{itemize}
\item Supplied local UnsolvedMath registry record, permanent ID \(30005445\), title \emph{Invariance Principles for Perturbed Balanced Environments}, exact source problem code \texttt{catalogue entry}. The record attributes the question to a 2023 Oberwolfach workshop report; no individual proposer is supplied.
\item \emph{Mini-Workshop: New Horizons in Motions in Random Media} (2023), Oberwolfach Reports, pp.\ 565--606, DOI \texttt{10.4171/owr/2023/10}. Supplied source report and literature-source link: \url{https://doi.org/10.4171/owr/2023/10}. No pinpoint page or contribution author for the quoted passage is supplied.
\item Supplied literature source, labelled \emph{random walks in random environments without ellipticity}: \url{https://doi.org/10.1016/j.spa.2013.01.007}. The supplied record does not identify a specific result or passage supporting either proposed extension.
\item Literature assessment dated exactly \texttt{2026-08-22}: the supplied curation labels the problem open, reports invariance principles for established balanced or perturbative regimes, and states that no proof covering either extension was found. This is supplied dated curation, not an independently verified assessment. The record itself requests independent verification of time-sensitive status and cited claims.
\item The supplied record labels its statement \texttt{corrected\_verified} and reports that the official report was checked directly and that the clean statement restores definitions and removes extraction artifacts. These are source metadata claims; the supplied excerpt still lacks the definitions needed for a precise formulation, and no independent source verification is asserted here.
\end{itemize}
\end{document}
